% Los preliminares no llevan cabecera ni numero, tampoco las paginas de relleno
\pagestyle{empty}

\chapter*{}

\thispagestyle{empty}

\begin{center}
{\large\bfseries \thetitle}\\
\end{center}
\begin{center}
\theauthor\\
\end{center}

\vspace{0.5cm}
\noindent{\textbf{Resumen}}\\

La soledad no deseada es un fenómeno que afecta a una parte considerable de la población pasando
desapercibida. En España alcanza al 20\% de los adultos y se agrava con la edad, hasta afectar a una de cada
cinco personas mayores de 75 años: se estima que cerca de tres millones de mayores de 65 la padecen, con un
coste para el Estado cifrado en torno a los 14.141 millones de euros anuales. Son numerosas las herramientas,
iniciativas y metodologías estudiadas para combatirla, algunas más fructíferas que otras.

El análisis del mercado actual revela, sin embargo, un hueco significativo. Los robots sociales de gama alta
ofrecen una experiencia cuidada a un precio inasumible para el pensionista medio español, mientras que los
compañeros conversacionales de uso masivo, optimizados para maximizar el tiempo de uso, no se diseñaron
pensando en personas mayores y han derivado en patrones documentados de dependencia emocional.

El objeto de este proyecto no pretende darle una solución a un problema de tal magnitud, sino contribuir a
mitigarlo desde el terreno que le es propio, el de la ingeniería del software. Se presenta el análisis, el
diseño y la implementación de \textbf{Quiet Companion}, una aplicación móvil que acompaña a la persona mayor
mediante un asistente conversacional por voz sobre hardware de uso común. La propuesta se articula
en torno a tres principios: la accesibilidad, con una interfaz que cumple los criterios de perceptibilidad y
operabilidad de la norma WCAG 2.2 y renuncia a cualquier gesto que no sea la pulsación simple; la revelación
progresiva, que reconoce que el mayor no es ni plenamente autónomo ni plenamente dependiente y le permite
gestionar la herramienta por sí mismo o delegarla en un familiar; y el cuidado distribuido, un modelo dual en
el que el familiar ajusta en remoto los parámetros de accesibilidad y los cambios se propagan en tiempo real
al dispositivo del mayor, de modo que la tecnología complemente el vínculo humano en lugar de sustituirlo.

El sistema se ha desarrollado en Flutter sobre una arquitectura limpia con organización por funcionalidades, y
se apoya en Firebase para la autenticación, la persistencia y la lógica de servidor, en un modelo de lenguaje
generativo y en servicios en la nube de reconocimiento y síntesis de voz. El resultado es un prototipo
funcional que cumple los objetivos comprometidos, con una limitación asumida que se propone como línea de
trabajo futuro: la ausencia de validación con usuarios reales que permita medir su efecto sobre la dimensión
social de la soledad.

\vspace{0.4cm}
\noindent{\textbf{Palabras clave}: Soledad no deseada, Personas mayores, Accesibilidad, WCAG 2.2, Asistente conversacional por voz, Inteligencia artificial generativa, Flutter, Firebase}\\

\chapter*{}

\thispagestyle{empty}

\begin{center}
{\large\bfseries \thetitleen}\\
\end{center}
\begin{center}
\theauthor\\
\end{center}

\vspace{0.5cm}
\noindent{\textbf{Abstract}}\\

Unwanted loneliness affects a considerable share of the population while going largely unnoticed. In Spain it
reaches 20\% of adults and worsens with age, up to one in five people over 75: close to three million people
over 65 are estimated to experience it, at a cost to the State of around 14.1 billion euros a year. Many
tools, initiatives and methodologies have been studied to counter it, some more fruitful than others.

A review of the current market reveals, however, a significant gap. High-end social robots offer a polished
experience at a price beyond the reach of the average Spanish pensioner, whereas mass-market conversational
companions, optimised to maximise time spent in the app, were not designed with older people in mind and have
led to documented patterns of emotional dependence.

This project does not claim to solve a problem of such magnitude, but to help mitigate it from the field it
belongs to, software engineering. It presents the analysis, design and implementation of \textbf{Quiet
Companion}, a mobile application that accompanies older people through a voice-based conversational assistant
on commodity hardware. The proposal rests on three principles: accessibility, through an interface
that meets the perceivability and operability criteria of the WCAG 2.2 standard and restricts interaction to a
single tap; progressive disclosure, which acknowledges that an older person is neither fully autonomous nor
fully dependent and lets them either manage the application themselves or delegate it to a relative; and
distributed care, a dual model in which the relative adjusts the accessibility settings remotely and the
changes propagate in real time to the older person's device, so that technology complements the human bond
instead of replacing it.

The system has been developed in Flutter on a clean architecture with feature-sliced organisation, and relies
on Firebase for authentication, persistence and server-side logic, on a generative language model, and on
cloud speech recognition and synthesis services. The result is a working prototype that meets the committed
objectives, with one acknowledged limitation proposed as future work: the absence of validation with real
users that would measure its effect on the social dimension of loneliness.

\vspace{0.4cm}
\noindent{\textbf{Keywords}: Unwanted loneliness, Older adults, Accessibility, WCAG 2.2, Voice conversational assistant, Generative artificial intelligence, Flutter, Firebase}\\

\chapter*{}
\thispagestyle{empty}

\noindent\rule[-1ex]{\textwidth}{2pt}\\[4.5ex]

D. \director, Profesor del Departamento Lenguajes y Sistemas de la Universidad de Granada. \\

D. \otrodirector, Profesor del Departamento Lenguajes y Sistemas de la Universidad de Granada. \\

\vspace{0.5cm}

\textbf{Informa:} \\

Que el presente trabajo, titulado \thetitle, ha sido realizado bajo su supervisión por \theauthor, y autorizamos la defensa de dicho trabajo ante el tribunal que corresponda.

\vspace{1cm}

Y para que conste, expiden el presente informe en \myLocation

\chapter*{}
\thispagestyle{empty}

\noindent\rule[-1ex]{\textwidth}{2pt}\\[4.5ex]

\textbf{Declaración de Originalidad} \\

Yo, \textbf{\theauthor}, alumno de la titulación \grado \ de la \myFacultyShort \ de la \myUni, con DNI \dni, declaro explícitamente que el trabajo presentado, correspondiente al curso académico \curso, es original, entendida esta, en el sentido de que no ha utilizado para la elaboración del trabajo fuentes sin citarlas debidamente.

\vspace{4cm}

\noindent Fdo: \theauthor

\vspace{1cm}

\begin{flushright}
\myLocation
\end{flushright}


\chapter*{}
\thispagestyle{empty}

\noindent\rule[-1ex]{\textwidth}{2pt}\\[4.5ex]

Yo, \textbf{\theauthor}, alumno de la titulación \grado \ de la \myFacultyShort \ de la \myUni, con DNI \dni, autorizo la ubicación de la siguiente copia de mi Trabajo Fin de Grado en la biblioteca del centro para que pueda ser consultada por las personas que lo deseen.

\vspace{4cm}

\noindent Fdo: \theauthor

\vspace{1cm}

\begin{flushright}
\myLocation
\end{flushright}

\vspace{3cm}

\doclicenseThis
El código fuente de este documento y de los programas desarrollados está disponible en el repositorio de
GitHub \url{https://github.com/maario106/quiet_companion}


\chapter*{Agradecimientos}
\thispagestyle{empty}

       \vspace{1cm}


Agradecimientos escritos...

% Se restituye el estilo de pagina para el resto del documento
\pagestyle{fancy}
